\scanpage{123}
\begin{flushright}
\fbox{\parbox{.53\linewidth}{Complexes de \uline{données} de De Rham\\ à p.d.}}
\end{flushright}
\begin{minipage}[t]{.43\linewidth}
Soient

$k$ anneau comm.

\deleted{\illegible{1}} \textcircled{$S$} une $k$-coalgèbre \inserted{ass. unitaire comm.}

\deleted{\illegible{1}}

\textcircled{$\bar S$} une sous-$k$-coalgèbre \inserted{ass. unit. comm.}
\[
\diagram{\begin{tikzcd}[row sep=large]
\bar S \arrow[d,Leftrightarrow]\\
S^+
\end{tikzcd}}
\]
\textcircled{$S^+$} $=S/\bar S$ le $k$-module quotient

[c'est une coalgèbre ass. comm. mais pas unitaire]

\textcircled{$\gamma_*$} un système de copuissances divisées sur $(S,S^+)$
\end{minipage}\hfill
\begin{minipage}[t]{.54\linewidth}
\[
\boxed{S\longrightarrow S\otimes_k S}
\]
[d'où \inserted{$V_k(S):$} $k'\longmapsto S^{k'}=\operatorname{Hom}_k(S,k')$ est un foncteur : \unsure{$\mathcal O_k$} dans les anneaux \unsure{commut.} \unsure{autrement dit} une $\mathcal O_k$-Alg. ass. uni comm.

[d'où \inserted{$V_k(\bar S):$} $k'\longmapsto\bar S^{k'}=\operatorname{Hom}_k(\bar S,k')$ est une $\mathcal O_k$-alg. ass. unit. comm. et $S^{k'}\longrightarrow\bar S^{k'}$ est un hom. de $k'$-Alg. ass. unit. comm. \inserted{i.e. $V_k(S)\longrightarrow V_k(\bar S)$ hom. de $\mathcal O_k$-algs.}

[$S^{+\,k'}\simeq\operatorname{Ker}(S^{k'}\longrightarrow\bar S^{k'})\subset S^{k'}$ est un idéal de $S^{k'}$, donc $V_k(S^+)\overset{\mathrm{dfn}}=V_k(S)^+$ est un Idéal de $V_k(S)$]

[d'où on a sur $V_k(\unsure{S^+|\bar S})=V_k(S)^+$ un syst. de puiss. divisées
\[
\gamma_i:V_k(S^+)\longrightarrow V_k(S^+)\qquad i\geq1
\]
définis par des hom de $k$-alg
\[
\gamma_i^*:\operatorname{Sym}_k(S^+)\longrightarrow\operatorname{Sym}_k^{\unsure{*}}(S^+)
\]
\deleted{\unsure{qui}} \unsure{des hom. de $k$-m}
\[
\gamma_i^*:S^+\longrightarrow\operatorname{Sym}_k^{\unsure{*}}(S^+)
\]
(qui, en vertu des axiomes sont en fait
\[
\boxed{\gamma_i^*:S^+\longrightarrow\operatorname{Sym}_k^i(S^+)\quad(i\geq1)}
\]
avec des axiomes dessus
\end{minipage}

\[
 t:S^+\longrightarrow k\quad\text{un $k$-hom.}\qquad
\text{i.e. }t\in V_k(S^+)(k)=V_k(S)^+(k)
\]
\figure{\begin{tikzpicture}[x=1cm,y=1cm]\draw (0,.65) .. controls (1,.85) and (.6,-.1) .. (1.6,-.1)--(2.05,-.1);\end{tikzpicture}}

Alors on définit un objet \unsure{simplicial} dans la cat. des foncteurs
\[
(k\text{-alg comm})\longrightarrow
\left(\begin{array}{l}\text{algèbres graduées alternées différentielles}\\
\text{(diff. de degré $+1$) à}\ \inserted{\text{avec idéaux à}}\\
\text{puiss. divisées (}\end{array}\right.
\]
et $\simeq$ dans la cat des \deleted{foncteurs} schémas affines \deleted{\illegible{1}} en algèbres graduées alternées différentielles (diff. de degré $+1$) avec Idéal \inserted{\uline{homogène}} à p. div. (qui est un schéma affine), les puiss. div. étant compatibles avec graduation et avec différentielle

\margin{122}
